\documentclass{ximera}
\title{Differentiating Trigonometric Functions using Limits}

\newcommand{\pskip}{\vskip 0.1 in}

\begin{document}
\begin{abstract}
Differentiating trig functions from the definition of the derivative.
\end{abstract}
\maketitle



\section{The Motion of a Balloon}

\begin{example} \label{ELdebar3}

Play the slider in Line 1 of the worksheet below to watch the motion of a balloon.

\begin{onlineOnly}
    \begin{center}
\desmos{xsk95weesg}{900}{600}
\end{center}
\end{onlineOnly}

\href{https://www.desmos.com/calculator/xsk95weesg}{151: Motion Cosine Derivative}

\begin{enumerate}
\item Sketch by hand a graph of the function
\[
    h  = f(t) \, , \, 0\leq t \leq 20
\]
that expresses the balloon's height (in feet) in terms of the number of seconds since the motion began. Label the axes with the appropriate variable names and units.

\item Activate the folder \emph{Graph} in Line 4 and play the motion again. How did you do?

\item Describe how the balloon's rate of ascent varies over the time of its motion. 

\item Are there time intervals during which the balloon's rate of ascent appears to be constant? If so, approximate these time intervals?

\item Use the motion to approximate the balloon's minimum and maximum rates of ascent. 

\item Sketch by hand a graph of the function
\[
     r = g(t) = \frac{dh}{dt} = \frac{d}{dt}\left( f(t) \right)
\]
that expresses the rate of ascent in terms of the number of seconds since the motion began. Label the axes with the appropriate variable names and units.
\end{enumerate}
\end{example}


\section{An Important Limit}
Computing the derivatives of the trigonometric functions hinges on one very important limit.

\begin{example}  \label{Ex8f7fd3lgzmq}

\begin{onlineOnly}
    \begin{center}
\desmos{xv3rn3tyxx}{900}{600}
\end{center}
\end{onlineOnly}

\href{https://www.desmos.com/calculator/xv3rn3tyxx}{151: Sin theta over theta Limit}

\begin{enumerate}
\item Find an expression for the function
\[
      r = m(\theta)
\]
that gives the average rate of change of the function
\[
     f(\theta) = \sin \theta
\]
over the interval $[0,\theta]$. Include the domain.

\item Drag the slider $u$ in Line 6 above to sketch by hand a graph of the function $m$. Label the axes with the appropriate varaible names and units.

\item Approximate the limit
\[
     \lim_{\theta \to 0} \frac{\sin\theta}{\theta}
\]
in two ways:

\begin{enumerate}
\item Visually, by zooming in on the origin in the worksheet above.

\item Numerically, with a table of average rates of change. Do this with a calculator and then open the \emph{average rates of change} folder in Line 4 as a check.
\end{enumerate}

\item What's your guess?
\[
   \lim_{\theta \to 0} \frac{\sin\theta}{\theta} = \answer{1} .
\]

\end{enumerate}
\end{example}



\end{document}